\documentclass[11pt,a4paper]{article}
\usepackage[utf8]{inputenc}
\usepackage[T1]{fontenc}
\usepackage[romanian]{babel}
\usepackage[margin=2cm]{geometry}
\usepackage{array}
\usepackage{tabularx}
\usepackage{enumitem}
\usepackage{graphicx}
\usepackage{setspace}
\usepackage{xcolor}
\usepackage{hyperref}
\usepackage{amssymb}
\setlength{\parindent}{0pt}
\newcommand{\fillline}[1][6cm]{\underline{\hspace{#1}}}
\newcommand{\fillbox}[1][0.4cm]{\framebox[#1]{\rule{0pt}{#1}}}
\newcommand{\checkboxempty}{$\square$}

\begin{document}

\begin{flushleft}
\textbf{PRIMĂRIA CLUJ-NAPOCA} \hfill Cod: 700.001 \hfill Nr. \fillline[4cm] / \fillline[2cm]
\end{flushleft}

\vspace{0.5cm}

\begin{center}
\textbf{CERERE PENTRU ELIBERAREA UNUI ABONAMENT DE PARCARE,\\
CONSTRUIT PRIN INVESTIŢIA STATULUI, PENTRU PERSOANE FIZICE}\\
\vspace{0.3cm}
\textbf{CĂTRE,\\
MUNICIPIUL CLUJ-NAPOCA\\
DIRECȚIA PARCĂRI PUBLICE}
\end{center}

\vspace{0.5cm}

Subsemnatul(a) \underline{ {{nume_complet}} } CNP \underline{ {{cnp}} } cu domiciliu/reședință la care se solicită locul de parcare în Cluj-Napoca, pe strada \underline{ {{zona_solicitata}} } nr. \fillline[1.5cm] bl. \fillline[1.5cm] ap. \fillline[1.5cm] telefon \fillline[4cm], solicit eliberarea unui abonament de parcare pentru autoturismul cu numărul de înmatriculare \underline{ {{nr_inmatriculare}} }, serie șasiu \underline{ {{marca_model}} }

\vspace{0.5cm}

Documente justificative după caz\textasciicircum{}:
\begin{itemize}[leftmargin=*]
\item[\checkboxempty] carte de identitate solicitant (domiciliul stabil sau reşedinţă în apropierea parcajului)
\item[\checkboxempty] certificat de înmatriculare a autoturism aflat în proprietatea solicitantului sau utilizatorului(față/verso) cu Inspecţia Tehnică Periodică valabilă-utilizatorul trebuie să fie menționat în talon.
\item[\checkboxempty] adeverință angajator în cazul autoturismelor de serviciu utilizate în mod permanent de către angajat;
\item[\checkboxempty] certificat de căsătorie în cazul în care autoturismul este bun comun al soților;
\item[\checkboxempty] certificat de încadrare în grad de handicap în cazul persoanelor cu dizabilități sau a reprezentanților legali ai acestora;
\end{itemize}

\vspace{0.3cm}

\begin{itemize}[leftmargin=*]
\item[\checkboxempty] Prin prezenta solicit comunicarea răspunsului pe următoarea adresă de e-mail:\\
\fillline[12cm]
\item[\checkboxempty] Mă oblig să comunic instituției orice modificare intervine în legătură cu această adresă de e-mail;
\item[\checkboxempty] Îmi exprim consimțământul ca Primăria Municipiului Cluj-Napoca să comunice orice informații, date personale, clarificări și completări pe adresa de e-mail indicată mai sus;
\item[\checkboxempty] Am luat la cunoștință faptul că în cazul nefuncționării serverului de e-mail comunicat sau în cazul adresei greșite de e-mail, Municipiul Cluj-Napoca nu poate fi tras la răspundere pentru acest lucru
\end{itemize}

\vspace{1cm}

DATA \fillline[5cm] \hfill SEMNĂTURA \fillline[5cm]

\vspace{0.5cm}

\textbf{* Cererea se va completa pe numele persoanei care figurează ca proprietar/coproprietar/utilizator în certificatul de înmatriculare al autoturismului în caz contrar nu va putea fi considerată validă.\\
** Locurile de parcare din cartierele Municipiului Cluj-Napoca se repartizează conform HCL nr. 412/2020, cu modificarile si completarile ulterioare, după data de 31 martie a anului în curs.}

\vspace{0.3cm}

\footnotesize
\textasciicircum{}Primăria municipiului Cluj-Napoca asigură gratuit fotocopierea documentelor.

Pentru documentele emise de către alte instituții publice, organe de specialitate ale administrației publice centrale și locale, precum și persoane juridice de drept privat, care potrivit legii au obținut statut de utilitate publică sau sunt autorizate să presteze un serviciu public, în regim de putere publică, pe care solicitantul nu le depune, acesta are posibilitatea de a-și da consimțământul expres ca Primăria municipiului Cluj-Napoca să obțină în numele său copii sau extrase ale acestora.

\vfill
\hrule
\vspace{0.2cm}
www.primariaclujnapoca.ro\\
telefon: 0264/596.030\\
Cluj-Napoca, str. Moţilor, nr. 3

\vspace{0.3cm}

Timp estimativ de completare: 5 minute

Datele personale care vă sunt solicitate prin prezenta cerere vor fi prelucrate numai în vederea procesării și soluționării solicitării dumneavoastră. Primăria municipiului Cluj-Napoca garantează securitatea procesării datelor și arhivarea acestora în conformitate cu prevederile legale în vigoare. Responsabilul Primăriei municipiului Cluj-Napoca cu protecția datelor poate fi contactat pe adresa de email dpo@primariaclujnapoca.ro.

În conformitate cu Regulamentul nr. 679/2016 aveți dreptul de a solicita Primăriei municipiului Cluj-Napoca, în ceea ce priveşte datele cu caracter personal referitoare la persoana vizată, accesul la acestea, rectificarea sau ştergerea acestora sau restricţionarea prelucrării sau a dreptului de a vă opune prelucrării, precum şi a dreptului la portabilitatea datelor.

\end{document}
